\section*{Capítulo \ref{ch:b1:photosynthesis} --- Fotossíntese e autotrofia}

\begin{solution}{exo:b1:photosynthesis:1}
As membranas externa e interna do envoltório, e a membrana do \omterm{def:b1:photosynthesis:chloroplast}{tilacoide}
por dentro. \omterm{prop:b1:photosynthesis:lightreactions}{Reações luminosas}: na membrana do \omterm{def:b1:photosynthesis:chloroplast}{tilacoide} (com prótons
bombeados para o lúmen). \omterm{def:b1:photosynthesis:fixation}{Ciclo de Calvin}: no \omterm{def:b1:photosynthesis:chloroplast}{estroma}.
\end{solution}

\begin{solution}{exo:b1:photosynthesis:2}
$2\,\mathrm{H_2O} + 2\,\mathrm{NADP^+} + 3\,\mathrm{ADP} + 3\,\mathrm{P_i}
\to \mathrm{O_2} + 2\,\mathrm{NADPH} + 2\,\mathrm{H^+} + 3\,\mathrm{ATP}$
com oito fótons. O oxigênio vem da água: algas que recebem
$\mathrm{H_2^{18}O}$ liberam $\mathrm{^{18}O_2}$; as que recebem
$\mathrm{C^{18}O_2}$ não (Ruben e Kamen).
\end{solution}

\begin{solution}{exo:b1:photosynthesis:3}
\omterm{def:b1:photosynthesis:fixation}{Fixação} (\omterm{def:b1:photosynthesis:fixation}{RuBisCO}), redução (a G3P), regeneração (do RuBP). Por
$\mathrm{CO_2}$: 3 \omterm{def:b1:water-small-molecules:atp}{ATP} e 2 NADPH.
\end{solution}

\begin{solution}{exo:b1:photosynthesis:4}
Perto de \qty{430}{nm} (azul) e de \qty{660}{nm} (vermelho). Em torno de
\qty{500}{nm} os \omterm{def:b1:photosynthesis:pigments}{carotenoides} absorvem e passam a energia à \omterm{def:b1:photosynthesis:pigments}{clorofila},
de modo que há liberação de oxigênio embora a própria \omterm{def:b1:photosynthesis:pigments}{clorofila} absorva
pouco ali.
\end{solution}

\begin{solution}{exo:b1:photosynthesis:5}
A luz vermelha distante (\qty{700}{nm}) sozinha move pouca fotossíntese;
acrescentada à luz de \qty{680}{nm}, ela dá mais que a soma das duas
velocidades. Dois sistemas de pigmentos com máximos de absorção
diferentes precisam cooperar em série, cada um precisando dos próprios
fótons: o \omterm{def:b1:photosynthesis:zscheme}{fotossistema} I (P700) e o \omterm{def:b1:photosynthesis:zscheme}{fotossistema} II (P680).
\end{solution}

\begin{solution}{exo:b1:photosynthesis:6}
$0.18\times 96\,500 = \qty{17.4}{kJ/mol}$ de prótons; $50/17.4 = 2.9$
prótons no mínimo; a sintase usa quatro, ou seja, \qty{72}{\%} de
eficiência.
\end{solution}

\begin{solution}{exo:b1:photosynthesis:7}
Prova que um gradiente de prótons através da membrana do \omterm{def:b1:photosynthesis:chloroplast}{tilacoide} basta
para mover a síntese de \omterm{def:b1:water-small-molecules:atp}{ATP}, sem luz e sem transporte de elétrons: o
acoplamento se dá pelo gradiente. Não prova, por si só, que as \omterm{prop:b1:photosynthesis:lightreactions}{reações
luminosas} fabricam \omterm{def:b1:water-small-molecules:atp}{ATP} desse modo na \omterm{def:b1:flowering-plant-organization:organs}{folha} (para isso são necessários o
desacoplador e as medidas de pH). Com um desacoplador o gradiente desaba
antes que a sintase consiga usá-lo: não há \omterm{def:b1:water-small-molecules:atp}{ATP}.
\end{solution}

\begin{solution}{exo:b1:photosynthesis:8}
$3\times 5 + 3\times 1 = 18$ carbonos; seis 3-PGA $= 18$; seis G3P
$= 18$; um G3P sai (3) e cinco G3P (15) regeneram três RuBP (15): $3 +
15 = 18$. Balanceado.
\end{solution}

\begin{solution}{exo:b1:photosynthesis:9}
Sem $\mathrm{CO_2}$ a \omterm{def:b1:photosynthesis:fixation}{RuBisCO} para: o RuBP deixa de ser consumido (sobe)
e o 3-PGA deixa de ser produzido (cai), enquanto a luz continua a
reduzi-lo. Sem luz não há \omterm{def:b1:water-small-molecules:atp}{ATP} nem NADPH: o 3-PGA deixa de ser reduzido
(sobe) e o RuBP deixa de ser regenerado (cai), enquanto a fixação
continua por pouco tempo.
\end{solution}

\begin{solution}{exo:b1:photosynthesis:10}
$E = hc/\lambda = 6.63\times 10^{-34}\times 3\times 10^8/6.8\times
10^{-7} = \qty{2.92e-19}{J}$; por mol, \qty{176}{kJ}. 48 fótons:
\qty{8450}{kJ}; eficiência $2870/8450 = \qty{34}{\%}$. Uma lavoura perde
a luz não absorvida (metade do espectro, reflexão, transmissão, falhas
entre as plantas), a fotorrespiração, a respiração das \omterm{def:b1:flowering-plant-organization:organs}{folhas}, dos
\omterm{def:b1:flowering-plant-organization:organs}{caules} e das raízes, a saturação luminosa das \omterm{def:b1:enzymes:enzyme}{enzimas} a pleno sol, e as
estações em que não há \omterm{def:b1:flowering-plant-organization:organs}{folhas}.
\end{solution}

\begin{solution}{exo:b1:photosynthesis:11}
Por 100 carboxilações, 33 oxigenações que liberam $16.7$ carbonos:
líquido de $83.3$, uma perda de \qty{17}{\%}. Custo de evitá-la: 2 \omterm{def:b1:water-small-molecules:atp}{ATP}
$\approx
0.2$ carbono por $\mathrm{CO_2}$, ou seja, \qty{20}{\%}; o projeto C4 só
se paga quando a perda supera cerca de um quinto --- em condições
quentes e secas, e não em condições frescas.
\end{solution}

\begin{solution}{exo:b1:photosynthesis:12}
Os dois usam membranas que bombeiam prótons com a energia do transporte
de elétrons e uma \omterm{thm:b1:photosynthesis:chemiosmosis}{ATP-sintase} que gasta o gradiente: o mecanismo é o
mesmo, e as sintases são homólogas. Mas os elétrons correm em sentidos
opostos --- da água ao NADPH no \omterm{def:b1:photosynthesis:chloroplast}{cloroplasto}, erguidos pela luz; do NADH
ao oxigênio na \omterm{def:b1:eukaryotic-cell:mitochondrion}{mitocôndria}, caindo
(\cref{ch:b1:respiration-fermentation}) --- e os ciclos do carbono não
são o inverso um do outro: o \omterm{def:b1:photosynthesis:fixation}{ciclo de Calvin} e o ciclo de Krebs não
compartilham nenhuma \omterm{def:b1:enzymes:enzyme}{enzima}, e a redução do $\mathrm{CO_2}$ usa NADPH ao
passo que a liberação dele usa $\mathrm{NAD^+}$. As equações globais são
inversas; a maquinaria é uma herança comum usada em dois sentidos, com
químicas diferentes na ponta do carbono.
\end{solution}

\begin{solution}{pb:b1:photosynthesis:1}
\textbf{1.} $hc/\lambda = \qty{2.92e-19}{J}$; $\times N_A = \qty{176}{kJ/mol}$.
\textbf{2.} \qty{700}{nm}: \qty{171}{kJ/mol}; \qty{450}{nm}:
\qty{266}{kJ/mol}. Um fóton azul excita a \omterm{def:b1:photosynthesis:pigments}{clorofila} a um estado mais
alto, que decai em picossegundos ao mesmo estado excitado mais baixo que
um fóton vermelho alcança; o excesso se perde como calor, e a química vê
um fóton nos dois casos.
\textbf{3.} $1/2.92\times 10^{-19} = \num{3.4e18}$ fótons por joule.
\textbf{4.} $\Delta E = 0.82 - (-0.32) = \qty{1.14}{V}$: $\Delta G =
96\,500\times 1.14 = \qty{110}{kJ}$ por mol de elétrons.
\textbf{5.} Dois fótons: cerca de $176 + 171 = \qty{347}{kJ}$;
armazenados, \qty{110}{kJ}: \qty{32}{\%}.
\textbf{6.} Quatro elétrons por $\mathrm{O_2}$ (duas moléculas de água),
dando dois NADPH: oito fótons.
\textbf{7.} 12 NADPH: 24 elétrons, 48 fótons.
\textbf{8.} $2.303\times 8.314\times 298\times 3/96\,500 = \qty{0.177}{V}$;
$\qty{17.1}{kJ}$ por mol de prótons.
\textbf{9.} $4 + 8 = 12$ prótons por $\mathrm{O_2}$: 3 \omterm{def:b1:water-small-molecules:atp}{ATP}.
\textbf{10.} $6\times 3 = 18$ \omterm{def:b1:water-small-molecules:atp}{ATP}: exatamente a necessidade do ciclo
nesse cálculo; na prática o rendimento fica mais perto de 2.6 por
$\mathrm{O_2}$ (\qty{4.7}{H^+} por \omterm{def:b1:water-small-molecules:atp}{ATP} no \omterm{def:b1:photosynthesis:chloroplast}{cloroplasto}) e o fluxo cíclico
em torno do \omterm{def:b1:photosynthesis:zscheme}{fotossistema} I fornece o restante.
\textbf{11.} $4\times 17.1 = \qty{68}{kJ}$ para um \omterm{def:b1:water-small-molecules:atp}{ATP} de \qty{50}{kJ}:
\qty{73}{\%}.
\textbf{12.} $10^{-5}\,\mathrm{mol/L}\times 10^{-15}\,\mathrm{L}\times
6\times 10^{23} = 6$ prótons livres no lúmen, contra $10^5$ por segundo
pelas sintases: o fluxo é carregado por prótons ligados aos grupos
tamponantes das \omterm{def:b1:proteins:peptide}{proteínas} e dos \omterm{def:b1:lipids:fattyacid}{lipídios} do lúmen, que os liberam tão
depressa quanto eles saem.
\textbf{13.} $18\times 50 + 12\times 220 = 900 + 2640 = \qty{3540}{kJ}$
para \qty{2870}{kJ} armazenados: \qty{670}{kJ} (um quinto) dissipados
como calor nas reações do ciclo, que precisam ser morro abaixo para
ocorrer.
\textbf{14.} $48\times 176 = \qty{8450}{kJ}$; $2870/8450 = \qty{34}{\%}$.
\textbf{15.} $0.34\times 0.45\times 0.9 = \qty{13.8}{\%}$.
\textbf{16.} $400\times 0.005\times 36\,000 = \qty{72}{kJ}$; glicose
$0.138\times 72\,000/2\,870\,000 = \qty{3.5e-3}{mol} = \qty{0.62}{g}$.
\textbf{17.} Líquido $0.6\times 0.62 = \qty{0.37}{g}$ de glicose,
$\qty{0.15}{g}$ de carbono ($72/180$).
\textbf{18.} A luz que cai entre as plantas ou no solo; as \omterm{def:b1:flowering-plant-organization:organs}{folhas} à
sombra de outras \omterm{def:b1:flowering-plant-organization:organs}{folhas} trabalhando abaixo da capacidade, e as \omterm{def:b1:flowering-plant-organization:organs}{folhas} a
pleno sol saturadas (as \omterm{def:b1:enzymes:enzyme}{enzimas} não conseguem usar todos os fótons); a
respiração das raízes, dos \omterm{def:b1:flowering-plant-organization:organs}{caules} e a noturna; a fotorrespiração; as
semanas antes de o dossel fechar e depois de ele senescer.
\textbf{19.} 33 oxigenações, 16.7 carbonos perdidos, 83.3 ganhos
líquidos.
\textbf{20.} Recuperação $33\times 3.5 = 117$ \omterm{def:b1:water-small-molecules:atp}{ATP}; ciclo 300 \omterm{def:b1:water-small-molecules:atp}{ATP}; total
de 417 \omterm{def:b1:water-small-molecules:atp}{ATP} para 83.3 carbonos: 5.0 \omterm{def:b1:water-small-molecules:atp}{ATP} por carbono líquido.
\textbf{21.} $3 + 2 = 5$ \omterm{def:b1:water-small-molecules:atp}{ATP} por carbono, sem nenhum carbono perdido ---
igual em \omterm{def:b1:water-small-molecules:atp}{ATP}, mas a planta C3 também perdeu \qty{17}{\%} da capacidade
de carboxilação dela; aproximadamente um empate a \qty{25}{\celsius}.
\textbf{22.} 50 oxigenações, 25 carbonos perdidos, 75 líquidos;
recuperação 175 \omterm{def:b1:water-small-molecules:atp}{ATP}; $475/75 = 6.3$ \omterm{def:b1:water-small-molecules:atp}{ATP} por carbono líquido, contra os 5
da \omterm{def:b1:photosynthesis:c4cam}{planta C4}, e com um quarto do carbono perdido: a C4 vence com folga.
\textbf{23.} Triplicar o $\mathrm{CO_2}$ desloca a competição na \omterm{def:b1:photosynthesis:fixation}{RuBisCO}
em favor da carboxilação, cortando a fotorrespiração e elevando a
fixação líquida nas plantas C3; as \omterm{def:b1:photosynthesis:c4cam}{plantas C4} já saturam a \omterm{def:b1:photosynthesis:fixation}{RuBisCO} com
$\mathrm{CO_2}$ concentrado e nada ganham além de uma pequena economia
de água.
\textbf{24.} \qty{1}{g} de carbono são \qty{0.083}{mol}: C3, $250\times 0.083 = \qty{21}{mol}$ de água, \qty{375}{g}; CAM, \qty{37}{g}.
\textbf{25.} Cerca de 8 fótons por $\mathrm{CO_2}$ (48 por glicose);
\qty{34}{\%} dos fótons absorvidos até a glicose; cerca de \qty{14}{\%}
da luz do sol incidente até o açúcar bruto e \qty{8}{\%} até o açúcar
líquido da \omterm{def:b1:flowering-plant-organization:organs}{folha} --- número que as perdas de uma planta inteira ao longo
de uma estação inteira reduzem a \qty{1}{\%}.
\end{solution}
